\section{Results and Discussion}
\label{sec:results}

\subsection{Nominal Performance at 25$^\circ$C and Relaxation Dynamics}
\label{sec:res_a}

Evaluated on the nominal $25^\circ\text{C}$ dataset with $70\%/30\%$ split, the 300-tree LightGBM teacher achieves MAE $= 0.488\%$ \cite{ke2017lightgbm}. Following the physics-calibrated distillation procedure, the 961-parameter student Micro-MLP records MAE $= 0.454\%$ and RMSE $= 1.236\%$ ($R^2 = 0.9986$), surpassing the teacher while using $>99.7\%$ fewer parameters. Smooth-$L_1$ Huber distillation filters tree-based step discontinuities, while physical descriptors ($\Delta V_{10}$, $\tau_{\text{log}}$, $P_{\text{ohm}}$) provide inductive bias preventing overfitting. As shown in Fig.~\ref{fig:nominal_and_rest}(a), predictions align tightly along the identity diagonal with zero-mean Gaussian residuals ($\sigma = 0.457\%$).

In EV applications, brief stops ($30$--$600\ \text{s}$) enable Coulomb counting drift reset \cite{che2025enhanced, waag2014critical}. Across rest horizons $t \in \{30, 60, 120, 300, 600\}\ \text{s}$, the student records MAEs of $1.170\%$, $0.787\%$, $0.523\%$, $0.339\%$, and $0.340\%$. As shown in Fig.~\ref{fig:nominal_and_rest}(b), direct $V_{\text{mea}}$ inversion fails catastrophically ($>10\%$ at $30\ \text{s}$), while the Stanford 3-stage RF remains above $2.0\%$ at $60\ \text{s}$ \cite{che2025enhanced}. By incorporating $\Delta V_{10}$ and $\tau_{\text{log}}$, the Micro-MLP back-extrapolates asymptotic diffusion recovery, enabling robust Coulomb counting re-initialization.

\begin{figure*}[!t]
	\centering
	\begin{subfigure}[b]{0.48\textwidth}
		\centering
		\includegraphics[width=\textwidth]{Figure_3_Parity_and_Error_Density_Ultimate_2.pdf}
		\caption{Parity correlation and residual density ($25^\circ\text{C}$)}
		\label{fig:parity_density}
	\end{subfigure}
	\hfill
	\begin{subfigure}[b]{0.48\textwidth}
		\centering
		\includegraphics[width=\textwidth]{Figure_4_Resting_Window_Sensitivity_Verified_2.pdf}
		\caption{Estimation MAE vs. rest horizon duration}
		\label{fig:rest_sensitivity}
	\end{subfigure}
	\caption{Nominal state estimation accuracy and resting horizon convergence. (a) Parity plot showing tight alignment ($R^2 = 0.9986$, $\text{MAE} = 0.454\%$, $\text{RMSE} = 1.236\%$) and zero-mean Gaussian residuals ($\sigma = 0.457\%$). (b) Error convergence across rest durations showing rapid stabilization below the $1.0\%$ commercial target within $60\ \text{s}$.}
	\label{fig:nominal_and_rest}
\end{figure*}

\subsection{Cross-Cell Generalization and Telemetry Noise Robustness}
\label{sec:res_c}

Following Stanford protocols \cite{che2025enhanced}, cells YX01--YX06 were used for training, while YX07 and YX08 served as unseen test targets under shifted profiles across $25^\circ\text{C}$ (Test 7) and $35^\circ\text{C}$ (Test 8). As depicted in Fig.~\ref{fig:gen_and_noise}(a), the student attains MAEs of $2.48\%$ (YX07) and $3.35\%$ (YX08), delivering $42.3\%$ and $22.1\%$ reductions over the Stanford limit ($4.30\%$) \cite{che2025enhanced}. Under multi-temperature stress, the model yields MAE $= 2.312\%$ ($\text{RMSE} = 4.420\%$) at $25^\circ\text{C}$ and $2.689\%$ ($\text{RMSE} = 5.376\%$) at $35^\circ\text{C}$. The minimal deviation ($\Delta \text{MAE} = 0.377\%$) confirms that Arrhenius scaling $\Phi_{\text{Arr}}$ stabilizes tracking across thermal envelopes.

In production BMS, voltage measurements are corrupted by inverter harmonics, EMI, and ADC quantization \cite{waag2014critical, wei2022future, rout2024online}. Under Gaussian noise $V_{\text{mea}}^{\text{noisy}} = V_{\text{mea}} + \delta v$ ($\delta v \sim \mathcal{N}(0, \sigma_v^2)$, $\sigma_v = 5.0\ \text{mV}$), the student maintains MAE $= 1.336\%$, outperforming the LightGBM teacher ($1.820\%$) and standard MLP ($3.240\%$) as shown in Fig.~\ref{fig:gen_and_noise}(b). Differential overpotential anchoring ($\Delta V_{10}$) cancels DC offsets, while Huber's linear regime for $|e| > \beta$ prevents noise spikes from distorting predictions.

\begin{figure*}[!t]
	\centering
	\begin{subfigure}[b]{0.48\textwidth}
		\centering
		\includegraphics[width=\textwidth]{Figure_5_Unseen_MultiTemp_Generalization_Camera_Ready_V2_2.pdf}
		\caption{Cross-cell and multi-temperature validation}
		\label{fig:unseen_multitemp}
	\end{subfigure}
	\hfill
	\begin{subfigure}[b]{0.48\textwidth}
		\centering
		\includegraphics[width=\textwidth]{Fig6_AFE_Noise_Robustness_2.pdf}
		\caption{AFE voltage noise resilience}
		\label{fig:afe_noise}
	\end{subfigure}
	\caption{Cross-cell generalization and sensor noise resilience. (a) Unseen cell validation on cells YX07 and YX08 benchmarked against the Stanford limit ($4.30\%$) across Test~7 ($25^\circ\text{C}$) and Test~8 ($35^\circ\text{C}$). (b) Estimation MAE progression under noise intensity $\sigma_v \in [0, 10]\ \text{mV}$ and error distribution under nominal $5.0\ \text{mV}$ perturbation.}
	\label{fig:gen_and_noise}
\end{figure*}

\subsection{Automotive Safety Compliance and Embedded Profiling}
\label{sec:res_f}

In ISO 26262 ASIL-D qualification \cite{yahyaei2023embedded}, $\ge 95\%$ of estimations must satisfy $|\Delta \text{SOC}| \le 2.0\%$. CDF analysis in Fig.~\ref{fig:comp_and_tradeoff}(a) confirms the student achieves $89.56\%$ within $\le 1.0\%$, $95.94\%$ within $\le 2.0\%$, and $97.56\%$ within $\le 3.0\%$, satisfying safety criteria.

The decisive criterion for production BMS deployment is sub-percent accuracy within microcontroller budgets ($M_{\text{Flash}} \le 16\ \text{kB}$, $T_{\text{exec}} \le 50\ \mu\text{s}$) \cite{yahyaei2023embedded}. On ARM Cortex-M4 at $160\ \text{MHz}$ (Table~\ref{tab:hardware_benchmark}), the Micro-MLP consumes $3.75\ \text{kB}$ Flash and $192\ \text{Bytes}$ SRAM with $12.0\ \mu\text{s}$ latency---a $154.2\times$ acceleration and $>99.7\%$ memory reduction over the Stanford RF baseline ($1420.0\ \text{kB}$, $1850.0\ \mu\text{s}$) \cite{che2025enhanced}, establishing embedded viability.

\begin{figure*}[!t]
	\centering
	\begin{subfigure}[b]{0.48\textwidth}
		\centering
		\includegraphics[width=\textwidth]{Figure_7_CDF_Authentic_IEEE_2.pdf}
		\caption{Empirical CDF and compliance breakdown}
		\label{fig:cdf_analysis}
	\end{subfigure}
	\hfill
	\begin{subfigure}[b]{0.48\textwidth}
		\centering
		\includegraphics[width=\textwidth]{Figure_8_Hardware_Accuracy_Tradeoff_Perfect_Final_230L_2.pdf}
		\caption{Accuracy-complexity trade-off on Cortex-M4}
		\label{fig:hardware_tradeoff}
	\end{subfigure}
	\caption{Automotive safety compliance and embedded execution profiling. (a) Cumulative error distribution verifying $95.94\%$ compliance within $|\Delta \text{SOC}| \le 2.0\%$. (b) Accuracy-complexity design space showing the Micro-MLP residing firmly within the ideal embedded domain (Flash $<16\ \text{kB}$, latency $<50\ \mu\text{s}$, MAE $<0.6\%$).}
	\label{fig:comp_and_tradeoff}
\end{figure*}

\subsection{Hardware Profiling, Latency, and Memory Footprint Benchmark}
\label{sec:res_g}

\begin{table}[!t]
	\caption{Hardware Benchmark on ARM Cortex-M4 @ 160~MHz}
	\label{tab:hardware_benchmark}
	\centering
	\resizebox{\columnwidth}{!}{%
		\begin{tabular}{lcccccc}
			\toprule
			\textbf{Model Architecture} & \textbf{Params} & \textbf{Flash} & \textbf{SRAM} & \textbf{Latency} & \textbf{Nominal MAE} & \textbf{ASIL-D} \\
			\midrule
			Stanford 3-Stage RF \cite{che2025enhanced} & $\approx 60\text{k}$ & $1420\ \text{kB}$ & $\approx 48\ \text{kB}$ & $1850.0\ \mu\text{s}$ & $1.630\%$ & No \\
			Tabular Transformer \cite{chen2024transformer} & $35.5\text{k}$ & $142.0\ \text{kB}$ & $\approx 18.5\ \text{kB}$ & $1250.0\ \mu\text{s}$ & $1.120\%$ & No \\
			Vanilla LSTM Baseline \cite{xiong2018towards} & $13.2\text{k}$ & $53.0\ \text{kB}$ & $\approx 6.4\ \text{kB}$ & $480.0\ \mu\text{s}$ & $0.940\%$ & No \\
			Standard MLP (No KD) & $3.2\text{k}$ & $12.8\ \text{kB}$ & $512\ \text{B}$ & $38.0\ \mu\text{s}$ & $0.980\%$ & Borderline \\
			PIML Teacher (LightGBM) \cite{ke2017lightgbm} & $\approx 18\text{k}$ & $285.0\ \text{kB}$ & $\approx 12.0\ \text{kB}$ & $720.0\ \mu\text{s}$ & $0.488\%$ & No \\
			\midrule
			\textbf{MicroPhys-BMS (Ours)} & \textbf{961} & \textbf{3.75 kB} & \textbf{192 B} & \textbf{12.0} $\boldsymbol{\mu}$\textbf{s} & \textbf{0.454\%} & \textbf{Yes} \\
			\bottomrule
		\end{tabular}%
	}
\end{table}

The student Micro-MLP was benchmarked against baselines on ARM Cortex-M4 at $160\ \text{MHz}$ (Table~\ref{tab:hardware_benchmark}). The Stanford RF baseline \cite{che2025enhanced} resides in the infeasible domain ($1420.0\ \text{kB}$ Flash, $1850.0\ \mu\text{s}$). Tabular Transformers and LSTM networks yield non-deterministic branching and high RAM, unsuitable for functional safety. In contrast, the proposed Micro-MLP is the only model in the \textit{Ideal Embedded Domain} (Flash $< 16\ \text{kB}$, latency $< 50\ \mu\text{s}$, MAE $< 0.60\%$), achieving $154.2\times$ latency acceleration and $>99.7\%$ memory reduction over ensemble baselines.
\subsection{Comparative Evaluation Against SOTA Baselines on Two Datasets}
\label{sec:res_h}

To assess the proposed \textit{MicroPhys-BMS} against SOTA paradigms, a comparative evaluation was conducted on two LFP datasets: Stanford ($25^\circ$C, 15{,}100 samples) \cite{che2025enhanced} and CALCE A123 ($24^\circ$C, 11{,}968 samples) \cite{calce2014a123, boyce2022calce}. Three benchmarks were implemented under identical protocols (70/30 split, same 9-D features): (i) PINN with 4 layers of 64 Tanh neurons \cite{karniadakis2021physics}; (ii) Tabular Transformer ($d_{\text{model}} = 27$, $n_{\text{head}} = 3$) \cite{chen2024transformer}; (iii) 5-model LightGBM ensemble (2{,}500 trees total) \cite{ke2017lightgbm}.

\subsubsection{Cross-Dataset Accuracy Comparison}
Table~\ref{tab:comparison_results} reports performance. On Stanford LFP, MicroPhys-BMS (961 params) achieves MAE $= 0.536\%$, competitive with LightGBM ($0.394\%$, $157{,}500$ params, $164\times$ more). PINN achieves $0.898\%$ ($13{,}185$ params), Transformer $1.188\%$ ($13{,}899$ params). On CALCE A123 (independent manufacturer, 18650 form), the proposed model achieves MAE $= 0.173\%$ with $98.89\%$ CDF within $|\Delta\text{SOC}| \le 2.0\%$, matching neural baselines with $13.7\times$--$164\times$ fewer parameters.

\begin{table}[!t]
	\caption{Cross-Dataset Comparative Evaluation}
	\label{tab:comparison_results}
	\centering
	\resizebox{\columnwidth}{!}{%
		\begin{tabular}{lcccc}
			\toprule
			\textbf{Model} & \textbf{Params} & \textbf{Dataset} & \textbf{MAE} & \textbf{RMSE} \\
			\midrule
			\multirow{2}{*}{MicroPhys-BMS (Ours)} & \multirow{2}{*}{\textbf{961}} & Stanford & $0.536\%$ & $1.492\%$ \\
			 & & CALCE A123 & $\mathbf{0.173\%}$ & $2.451\%$ \\
			\midrule
			\multirow{2}{*}{PINN \cite{karniadakis2021physics}} & \multirow{2}{*}{$13{,}185$} & Stanford & $0.898\%$ & $2.050\%$ \\
			 & & CALCE A123 & $0.137\%$ & $2.477\%$ \\
			\midrule
			\multirow{2}{*}{Transformer \cite{chen2024transformer}} & \multirow{2}{*}{$13{,}899$} & Stanford & $1.188\%$ & $2.367\%$ \\
			 & & CALCE A123 & $0.132\%$ & $2.484\%$ \\
			\midrule
			\multirow{2}{*}{LightGBM Ens. \cite{ke2017lightgbm}} & \multirow{2}{*}{$\approx 157{,}500$} & Stanford & $\mathbf{0.394\%}$ & $\mathbf{1.037\%}$ \\
			 & & CALCE A123 & $0.231\%$ & $1.299\%$ \\
			\bottomrule
		\end{tabular}%
	}
\end{table}

\subsubsection{Noise Robustness and CDF Analysis}
As depicted in Fig.~\ref{fig:cdf_comparison}(a), MicroPhys-BMS achieves $93.83\%$ CDF within $|\Delta\text{SOC}| \le 2.0\%$ on Stanford, outperforming PINN ($85.61\%$) and Transformer ($79.33\%$). On CALCE A123, all models achieve $>97.5\%$. Under $\pm 5.0$~mV noise, the proposed model maintains MAE $= 1.343\%$ (Stanford) and $0.249\%$ (CALCE), both within the $1.456\%$ target. As shown in Fig.~\ref{fig:cdf_comparison}(b), the 961-parameter Micro-MLP is the only model satisfying embedded constraints (Flash $< 16$~kB, ASIL-D) while maintaining competitive accuracy, whereas LightGBM requires $157{,}500$ parameters.

\begin{figure*}[!t]
	\centering
	\begin{subfigure}[b]{0.48\textwidth}
		\centering
		\includegraphics[width=0.85\textwidth]{Figure_7_CDF_Comparison.pdf}
		\caption{CDF of absolute SOC error}
		\label{fig:cdf_comparison_a}
	\end{subfigure}
	\hfill
	\begin{subfigure}[b]{0.48\textwidth}
		\centering
		\includegraphics[width=0.85\textwidth]{Figure_8_Hardware_Tradeoff_Comparison.pdf}
		\caption{Accuracy-complexity trade-off}
		\label{fig:hardware_tradeoff_comparison}
	\end{subfigure}
	\caption{Comparative evaluation. (a) CDFs on both datasets: MicroPhys-BMS achieves $93.83\%$ (Stanford) and $98.89\%$ (CALCE) within $|\Delta\text{SOC}| \le 2.0\%$. (b) Accuracy-complexity: the 961-param Micro-MLP resides in the ideal embedded domain (Flash $< 16$~kB), while baselines need $13{,}185$--$157{,}500$ params.}
	\label{fig:cdf_comparison}
\end{figure*}

